\documentclass{situationist-preprint}

\title{What Remains of Aquinas's Doctrine}
\msid{REL-16-b}
\msfamily{REL-16}
\msdiscipline{Philosophy of Religion}
\mstype{Companion manuscript}
\mskeywords{divine simplicity; Aquinas; aseity; knowledge of vision; accidents; contingent knowledge}

\begin{document}
\maketitle

\begin{abstract}
REL-16-a argues that divine simplicity is incoherent with God's knowledge of contingent truths, using only the thesis that God's knowledge is fixed by God's essence. Aquinas's doctrine is a bundle of distinct denials of composition, and the argument does not touch all of them. We separate five theses in \emph{Summa Theologiae} I, q.~3, and ask which survive. God's immateriality, the identity of essence and existence, and the absence of genus and difference survive intact; the mutual identity of God's necessary attributes survives in a restricted form. What must go is the denial of accidents, at least for God's knowledge of what is contingent, and with it the identity of that knowledge with the essence.
\end{abstract}

\section{Five theses}

Aquinas argues in \emph{Summa Theologiae} I, q.~3, that God is in no way composite, by denying in turn each mode of composition he recognizes \citep{AquinasST}. We separate the theses as follows.
\begin{enumerate}
\item[(T1)] God is not a body and is not composed of matter and form (aa.~1--2).
\item[(T2)] God is identical with God's own essence or nature (a.~3).
\item[(T3)] In God essence and existence are the same (a.~4).
\item[(T4)] God is not in a genus and is not composed of genus and difference (a.~5).
\item[(T5)] There are no accidents in God (a.~6); hence God's attributes, knowledge included, are identical with one another and with the essence.
\end{enumerate}
The argument of REL-16-a uses Definition 2.1, that God's knowledge is fixed by God's essence, with Assumptions 2.2--2.5: God's essence is the same at every world, knowledge is factive, God knows every truth, and some truths are contingent. Theorem 3.2 concludes that God's knowledge cannot be fixed by the essence. Of the five theses, only (T5) entails Definition 2.1. The others say nothing about what fixes God's knowledge, and the argument leaves them standing.

\begin{proposition}\label{p:survive}
(T1)--(T4) are consistent with the conclusion of REL-16-a.
\end{proposition}

\begin{proof}
None of (T1)--(T4) entails that God's knowledge at a world is a function of God's essence. (T1) and (T4) concern what God is not composed of; (T2) and (T3) concern the relation of God to God's essence and of essence to existence. A being satisfying all four could have, in addition to its necessary essence, an intrinsic state of knowledge that differs between worlds.
\end{proof}

(T3) is the heart of the doctrine for Aquinas, and it is the thesis that secures aseity: a being whose essence is its existence depends on nothing for its being. It survives.

\section{The restricted identity of attributes}

Aquinas distinguishes God's knowledge of simple intelligence, of everything possible, from the knowledge of vision, of what is actual at some time (I, q.~14, a.~9). The first does not vary between worlds: what is possible is necessarily possible. It can therefore be fixed by the essence without contradiction.

\begin{proposition}\label{p:restricted}
The identity of God's necessary attributes with one another and with the essence, including the knowledge of simple intelligence, is consistent with REL-16-a.
\end{proposition}

\begin{proof}
The proof of Theorem 3.2 of REL-16-a requires that the knowledge fixed by the essence include some proposition true at one world and false at another (Assumption 2.5). The knowledge of simple intelligence includes no such proposition, since its objects are necessary truths about what is possible. Fixing it by the essence yields no contradiction.
\end{proof}

What does not survive is the knowledge of vision. God knows that Socrates sits at $t$, and Aquinas holds that God wills things other than himself without necessity (I, q.~19, a.~3); at worlds where God wills otherwise, God knows otherwise. That knowledge is intrinsic to God, as REL-16-a argues against the extrinsic reply, and it is not fixed by the essence. It is an accident in Aquinas's sense, and (T5) fails.

\begin{objection}
The denial of accidents is required by (T3): a being whose essence is existence is pure act, and an accident is a potency actualized.
\end{objection}

\begin{reply}
An accident of the knowledge of vision is not a potency actualized from outside. Its content depends on what God wills, and, where creatures are free, on what they choose, which REL-05 concludes God's knowledge depends on counterfactually. Dependence of content is not dependence of being. \citet{StumpKretzmann1985} defend simplicity by treating such knowledge as an extrinsic relation; once that reply is rejected, the defender of (T3) must accept that pure act is compatible with contingent intrinsic states, or else that God does not know what is contingent. The objection of \citet{Plantinga1980}, that the identity theses make God a property, and the truthmaker reading of those theses in \citet{Brower2008}, both concern (T2) and (T5) as theses about attributes in general; the present point is narrower and holds on either reading.
\end{reply}

\section{Conclusion}

Simplicity must be abandoned as a doctrine about all of God's attributes, but (T1)--(T4) and the unity of God's necessary attributes survive. What goes is the thesis that God has no accidents, and with it the claim that God's knowledge of what happens is the same as what God is.

\begin{thebibliography}{99}\small
\bibitem[Aquinas(1265--1274)]{AquinasST} Thomas Aquinas (1265--1274). \emph{Summa Theologiae}, Prima Pars.
\bibitem[Brower(2008)]{Brower2008} Brower, J.~E. (2008). Making sense of divine simplicity. \emph{Faith and Philosophy} 25(1): 3--30.
\bibitem[Plantinga(1980)]{Plantinga1980} Plantinga, A. (1980). \emph{Does God Have a Nature?} Marquette University Press.
\bibitem[Stump and Kretzmann(1985)]{StumpKretzmann1985} Stump, E. and Kretzmann, N. (1985). Absolute simplicity. \emph{Faith and Philosophy} 2(4): 353--382.
\end{thebibliography}

\end{document}
